\section{Lecture 17: Activation Functions and Depth}

\subsection{Motivation}

In the previous lecture, we saw that training deep networks is difficult due to unstable gradient flow.

\medskip

\noindent
\textbf{Key question.}  
How can we design networks that both:
\begin{itemize}
    \item are expressive,
    \item and are trainable?
\end{itemize}

\medskip

\noindent
\textbf{Guiding idea.}  
Activation functions and architectural design play a central role in controlling both expressivity and gradient flow.

\subsection{Nonlinear Activations}

A neural network layer has the form:
\[
h = \sigma(Wx + b),
\]
where $\sigma$ is a nonlinear activation function.

\medskip

\noindent
\textbf{Why nonlinearity is necessary.}

\begin{itemize}
    \item Without $\sigma$, composition of layers is linear:
    \[
    W_L \cdots W_1 x = W x
    \]
    \item Depth would not increase expressivity
\end{itemize}

\medskip

\noindent
\textbf{Common activation functions.}

\begin{itemize}
    \item \textbf{Sigmoid:}
    \[
    \sigma(z) = \frac{1}{1 + e^{-z}}
    \]

    \item \textbf{Tanh:}
    \[
    \tanh(z)
    \]

    \item \textbf{ReLU:}
    \[
    \text{ReLU}(z) = \max(0, z)
    \]
\end{itemize}

\medskip

\noindent
\textbf{Comparison.}

\begin{itemize}
    \item Sigmoid / tanh:
    \begin{itemize}
        \item smooth
        \item saturate $\Rightarrow$ vanishing gradients
    \end{itemize}

    \item ReLU:
    \begin{itemize}
        \item non-saturating for $z > 0$
        \item preserves gradient magnitude
    \end{itemize}
\end{itemize}

\medskip

\noindent
\textbf{Insight.}  
Activation functions control both expressivity and gradient flow.

\subsection{Architectural Design Patterns}

Beyond individual layers, architecture matters.

\medskip

\noindent
\textbf{Basic pattern.}
\begin{itemize}
    \item Linear transformation $\rightarrow$ nonlinearity
\end{itemize}

\medskip

\noindent
\textbf{Stacking layers.}
\begin{itemize}
    \item Builds hierarchical representations
\end{itemize}

\medskip

\noindent
\textbf{Key considerations.}

\begin{itemize}
    \item Depth (number of layers)
    \item Width (number of units per layer)
    \item Activation choice
\end{itemize}

\medskip

\noindent
\textbf{Design principle.}  
Structure should reflect the underlying structure of the data.

\subsection{Depth versus Width}

\textbf{Two ways to increase model capacity:}

\begin{itemize}
    \item Increase width (more neurons per layer)
    \item Increase depth (more layers)
\end{itemize}

\medskip

\noindent
\textbf{Width.}
\begin{itemize}
    \item Can approximate complex functions (universal approximation)
    \item May require exponentially many units
\end{itemize}

\medskip

\noindent
\textbf{Depth.}
\begin{itemize}
    \item Enables compositional structure
    \item More parameter-efficient
\end{itemize}

\medskip

\noindent
\textbf{Key intuition.}

\begin{itemize}
    \item Depth allows reuse of intermediate features
    \item Represents hierarchical computations
\end{itemize}

\medskip

\noindent
\textbf{Tradeoff.}
\begin{itemize}
    \item Deeper networks are harder to train
\end{itemize}

\medskip

\noindent
\textbf{Insight.}  
Depth provides efficient expressivity but introduces optimization challenges.

\subsection{Expressivity Considerations}

\textbf{Piecewise linear perspective (ReLU networks).}

\begin{itemize}
    \item ReLU networks partition input space into regions
    \item Each region corresponds to a linear function
\end{itemize}

\medskip

\noindent
\textbf{Effect of depth.}

\begin{itemize}
    \item Number of linear regions grows rapidly with depth
    \item Enables highly complex decision boundaries
\end{itemize}

\medskip

\noindent
\textbf{Representation viewpoint.}

\begin{itemize}
    \item Early layers: simple features
    \item Deeper layers: abstract features
\end{itemize}

\medskip

\noindent
\textbf{Connection to kernels.}

\begin{itemize}
    \item Kernel methods use fixed feature maps
    \item Deep networks learn feature maps dynamically
\end{itemize}

\medskip

\noindent
\textbf{Insight.}  
Expressivity depends on both architecture and learned representation.

\subsection{A Unifying Perspective}

Activation functions and depth together determine:

\begin{itemize}
    \item \textbf{Expressivity:} ability to represent complex functions
    \item \textbf{Trainability:} stability of gradient flow
    \item \textbf{Inductive bias:} structure imposed on solutions
\end{itemize}

\medskip

\noindent
\textbf{Core idea.}  
Deep networks balance expressivity and trainability through architectural design.

\medskip

\noindent
\textbf{Key insight.}  
The success of deep learning depends as much on architecture as on optimization.

\subsection{Outlook}

In this lecture, we studied:

\begin{itemize}
    \item nonlinear activation functions,
    \item architectural design patterns,
    \item depth versus width tradeoffs,
    \item expressivity of neural networks.
\end{itemize}

\medskip

\noindent
However, activation functions alone do not fully solve training instability.

\medskip

\noindent
In the next lecture, we study normalization methods, which stabilize training and enable deeper networks.

\medskip

\noindent
\textbf{Guiding principle.}  
Designing deep networks requires balancing expressivity with stable gradient propagation.